\scanpage{33}
% Pencil archive foliation: 32.
on a mieux :
\[
M\in\operatorname{Ob}\mathcal{M}_i^+(X)\Longleftrightarrow
M(-j)\in\operatorname{Ob}\mathcal{M}^{+i+2j}(X)
\]
De cette façon, la filtration \inserted{de $\mathcal{M}^+(X)$} par les
$\mathcal{M}_i^+(X)$ \inserted{($i\geq0$)} peut se prolonger en une filtration
de $\mathcal{M}(X)$ par des $\mathcal{M}^i(X)$ ($i\in\mathbb{Z}$) de
telle façon que pour $i\geq0$, on ait la filtration
déjà \deleted{\unsure{\mathcal{M}^+(X)}} envisagée, et pour $i$ quelconque,
on définit
\[
M\in\operatorname{Ob}\mathcal{M}^i(X)\Longleftrightarrow
M(-j)\in\operatorname{Ob}\mathcal{M}^{i+2j}(X)
\]
si on prend $j$ assez grand pour que
$M(-j)\in\operatorname{Ob}\mathcal{M}^+(X)$. On notera que
\[
\mathcal{M}^{+i}(X)=\mathcal{M}^i(X)\cap \mathcal{M}^+(X)
\]
\margin{indices en bas}
si $i\geq0$, et pour tt $i$ si on définit
\[
\mathcal{M}^{+i}(X)=\{0\}\deleted{\unsure{\mathcal{M}^+(X)}}\quad\text{si }i<0
\]
\deleted{(\unsure{on voit} par là \illegible{3} que $\mathcal{M}^{+i}(X)$
\inserted{\deleted{$\mathcal{M}^+(X)$}}
\unsure{croisse avec $i$})}
On aura
\[
\otimes\unsure{\mathbb{Q}|\mathbb{Z}}(-1):
\mathcal{M}^i(X)\overset{\text{équiv}}{\simeq}\mathcal{M}^{i+2}(X)
\]
\[
\otimes\unsure{\mathbb{Q}|\mathbb{Z}}(j):
\mathcal{M}^i(X)\overset{\text{équiv}}{\simeq}\mathcal{M}^{i-2j}(X)
\]
